% ============================================================
% Chapter 22 — Demonstration Plan
% Source: PSYCON Engineering Design Document, Vol. 4, Ch. 22
% ============================================================
\section{Demonstration Plan}
\label{sec:demonstration-plan}

A live demonstration should be prepared, showing:
\begin{itemize}
    \item Device startup.
    \item Sensor initialization.
    \item Live physiological acquisition.
    \item Live audio capture (with consent).
    \item Data synchronization.
    \item Backend visualization.
    \item AI inference output.
    \item Consented wearer enrollment and encrypted-profile deletion.
    \item Anonymous speaker timeline, conservative wearer match, and attributed transcript.
    \item Per-speaker speaking rates, pauses, response gaps, overlaps, interruptions, and jitter with unavailable reasons.
    \item Safe shutdown.
\end{itemize}

\begin{requirementbox}
A recorded backup demonstration should also be prepared in case of hardware issues.
\end{requirementbox}

The Week~3 software demonstration uses a consented multilingual recording with at least two speakers, natural pauses, and overlap. It must show that missing model access, missing enrollment, ambiguous identity, and poor-quality speech produce explicit abstention states while acoustic analysis and transcription remain independently usable.
